%%
%% CONCLUSION
%%
\chapter{Conclusion}

Nous avons expos\'e les raisons pour lesquelles nous d\'efendons une
approche bottom-up g\'en\'erale en vue de la r\'ealisation d'un
environnement de d\'eveloppement cellulaire \`a vocation universelle. %
L'ampleur de la t\^ache, qui reste largement inachev\'ee, mais aussi
l'exp\'erience acquise, nous conduit \`a pr\'eciser quelques directions
privil\'egi\'ees de d\'eveloppement ult\'erieur.

Les principaux r\'esultats portent sur l'obtention d'un moteur
cellulaire rapide et sur la validation de plusieurs outils bien
adapt\'es \`a l'analyse et la pr\'esentation visuelle des dynamiques
g\'en\'er\'ees par le Workbench.

La r\'ealisation des m\'ecanismes de sp\'ecification des architectures
cellulaires pose un probl\`eme de bootstrap. %
Ce probl\`eme se pose \'egalement pour la r\'ealisation d'une biblioth\`eque
d'objets cellulaires \'el\'ementaires. %
Ainsi, nous pensons qu'il convient de consolider le noyau de
fonctionnalit\'es acquises, par une int\'egration plus \'etroite de quelques
moteurs cellulaires simples et rapides aux environnements externes,
tout en poursuivant la recherche des m\'ethodes utilisables pour la
d\'efinition d'architectures plus g\'en\'erales. %
N\'eanmoins, la r\'ealisation d'outils puissants~(rapides et flexibles) de
manipulation de tables de transitions, n'\'eloignerait pas la tentation
d'exp\'erimenter des AC moins radicalement bottom-up.

%% \section{Un AC d'un autre type}

\mydumpfull[t]{xu-photo-1}%
{}
{La dynamique de l'AC \arule{xu}}%
{Un cube d'espace-temps d'ar\^ete~128, rendu par ray-tracing
  d'iso-sur\-face. %
  De tels traitements sont, avec les outils que nous avons utilis\'es \`a
  la limite des capacit\'es des ressources dont nous disposions.}

\mydumpfull[t]{xu-photo-1-top}%
{}%
{Un rendu d'une tranche des valeurs du dernier
  plan de la figure~\protect\mygetdump{xu-photo-1}}%
{La photo~\mygetdump{xu-photo-1} permet de ce faire une id\'ee de la
  dynamique de croissance de la r\`egle, mais elle ne permet pas de
  distinguer le pattern de distribution \'energ\'etique de chaque tranche
  de temps.}

Nous pr\'esentons donc en conclusion un AC dont la r\'ealisation effective
exclut l'usage d'une combinatoire de tables de transitions. %
Cet AC utilise un voisinage de von Neumann, l'horloge interne est \`a
deux pas, et la structure cellulaire est \`a six champs. %
Chaque champ est un entier, dont la taille est, dans notre
r\'ealisation, de 32~bits. %
La fonction de transition est exprim\'ee\footnote{%
  Contrairement \`a toutes les autres r\`egles sp\'ecifi\'ees dans la pr\'esente \'etude.}%
en terme de modification simultan\'ee des cinq cellules du voisinage,
avec consultation de l'\'etat local uniquement. %
La r\`egle de transition r\'ealise une \aquote{physique} \'energ\'etique
qualitative. %
Les quatre premiers champs de la variable d'\'etat cellulaire associent
des poids aux quatre directions de l'espace~(nord, sud, est et
ouest). %
Les deux derniers champs associent des potentiels cin\'etiques aux axes
\bialt{nord}{sud} et \bialt{est}{ouest}. %
Les poids sont des donn\'ees initiales. %
La somme des poids sur l'ensemble des cellules de l'espace reste
invariante par conservation de l'\'energie. %
Les potentiels associ\'es aux axes~NS et~EO sont fonction du cumul des
poids des deux directions de chaque axe. %
Le rapport de cette quantit\'e \`a la somme des poids des quatre
directions d\'etermine le d\'eplacement additif, champ par champ, de la
cellule vers une de ses quatre voisines par fusion. %
Ce mouvement-fusion a lieu au temps~2 de chaque cycle. %
L'op\'eration de fusion porte \'egalement sur les potentiels, dont les
reliquats sont conserv\'es. %

Par exemple pour l'axe EO:

\footnotesize
\begin{verbatim}
        /* Code du choix de direction du mouvement */
        dx = icell->poids[EST] - icell->poids[OUEST];
        icell->potentiel[EO] += dx;

        if (abs(icell->potentiel[EO]) >= sum) {
                if (dx > 0) {
                        icell->potentiel[EO] -= sum;
                        x = modx(x + 1);
                } else {
                        icell->potentiel[EO] += sum;
                        x = modx(x - 1);
                }
        }
\end{verbatim}
\normalsize

Au temps~1 de chaque cycle, la somme est compar\'ee \`a la constante
universelle~K de notre physique. %
Si la somme est sup\'erieure \`a~K, la cellule explose. %
Cette explosion distribue la cellule dans les quatre voisins par
l'op\'eration de fusion. %
La division de la cellule en quatre parties est une division enti\`ere:
le reliquat de cette division permet d'ajouter la rupture de sym\'etrie
\`a la conservation de la quantit\'e allou\'ee au temps z\'ero. %
Le choix de la distribution du reliquat de cette division constitue
avec le choix de l'espace initial et de la valeur de~K les seuls
param\`etres contr\^olables. %
Un reliquat subit une rotation et une translation relative \`a sa
position cellulaire. %
Les exemples de dynamique expos\'es le sont
pour K = 4, rotation = clockwise next, translation = clockwise next. %

la figure~\mygetdump{xu-photo-1} est une photo de l'\'evolution~(de bas en
haut) de l'espace initial suivant:

\footnotesize
\begin{verbatim}
X       Y       N       S       E       O
-60     0       0       100     100000  0
60      8       0       100     100     100000
0       -60     0       100000  0       0
0       48      100000  0       200     400
\end{verbatim}
\normalsize

de taille 128 par 128, sur 128 g\'en\'erations. %
L'origine du rep\`ere est au centre.

\vspace{4cm}
Les r\'ealisations et propositions que nous avons expos\'ees
constituent des \'el\'ements d'une \'etude pratique des espaces du calcul
cellulaire reposant sur un usage g\'en\'eralis\'e de tables de transitions
qui permet la construction de moteurs rapides pour le calcul
cellulaire sur machines s\'equentielles. %
Nous avons montr\'e comment cette approche permet des transformations
cellulaire-SIMD. %
Nous avons \'egalement propos\'e l'usage d'un nouveau param\`etre de contr\^ole
sur les espaces de r\`egles et des techniques de visualisation
sp\'eciales pour les espaces-temps cellulaire. %
Nous comptons utiliser ces \'el\'ements pour base de nouveaux
d\'eveloppements int\'egrant la dimension esth\'etique des visualisations
comme m\'ethode d'\'evaluation heuristique des propri\'et\'es morphog\'en\'etiques
des r\'egions d'espaces du calcul cellulaire visit\'ees.

\mydumpfull{xu-Cfilm-1-1}%
{}%
{D'autre points de vue}%
{La pr\'esentation d'une photo d'une projection~2D d'un plongement~3D de la
  dynamique d'un r\`egle sur un espace cellulaire, ne remplace pas
  l'exp\'erience d'un pr\'esentation incarn\'ee dans le temps\footnote{%
    A priori un film classique, respectant le temps propre de la
    dynamique, mais sans exclure la possibilit\'e que des distortions de
    l'axe de projection temporelle soient porteuse d'informations
    pertinentes.}, %
  qui plus est, ces photos appellent par elle-m\^emes le besoin d'un
  film de la variation des param\`etres de projections.}

\mydumpfull{xu-Cfilm-1-2}%
{}%
{Et un \'epluchage}%
{R\'ealis\'e en variant lin\'eairement la valeur de seuil d'extraction d'iso-sur\-face}

\mydumpfull{xu-Cfilm-2-1}%
{}%
{Une autre configuration initiale}%
{}

\mydumpfull{xu-Cfilm-2-2}%
{}%
{Et un \'epluchage exponentiel}%
{}

%%
%% Local Variables:
%% mode: latex
%% mode: auto-fill
%% iso-accents-minor-mode: nil
%% TeX-master: "top"
%% TeX-command-default: "mytex"
%% Local IspellParsing: "latex-mode ~latin1"
%% Local IspellPersDict: "~/.ispell_lisp"
%% LocalWords: "bottom-up Workbench NS EO dx icell icell abs sum modx else next"
%% LocalWords: "clockwise clockwise next xu-photo l'AC xu xu-Cfilm bottom"
%% LocalWords: "Neumann Cfilm" cellulaire-SIMD espaces-temps
%% End:
